% ============================================
% DOC-01 (ZH)
% 中文翻译 — SPEC-935-R356
% ============================================
\section*{\textbf{DOC-01} --- 掘尸}
\markboth{DOC-01 --- 掘尸}{}

\noindent \textit{掘墓人热内西奥·巴尔博萨的记录}

\noindent \textbf{掘墓人热内西奥·巴尔博萨的记录}
\textit{1979年6月14日向收容医院领导层所作证词的逐字转录。文件按C-1260-79号归档。}

\vspace{0.5em}\noindent\rule{\textwidth}{0.5pt}\vspace{0.5em}

\noindent 那是当天挖的第三座坟。我们认得老坟，因为泥土已经沉实了。新坟的土是松的，几乎用手就能刨开。老坟不一样。土变硬了，上面结了一层壳，你一锹下去，像是要碰到石头。

\noindent 1,260的坟是老的。非常老。按地图算，1917年。最早埋在那里的那批人之一。

\noindent 我挖。挖了大约四十厘米，锹碰到了木头。再挖。棺材——是个棺材，但又不是。你知道吗，木头烂掉就变成了土？不是消失。它变成了别的东西。一摊深色的糊，还保留着原来形状。

\noindent 我用锹撬开。不是不敬。是拉不开盖子——它已经是地面的一部分了。

\noindent 我在里面看见的不是一具骷髅。

\noindent 是一具尸体。

\noindent 不，我没疯。是一具尸体。皮肤是深色的，几乎发黑，但\textit{完好}。像是上周才埋的。您可以问尼洛先生，他当时跟我在一起。他也看见了。脸上的皮肤还包着骨头。嘴张着。眼睛——

\noindent 眼睛不在了。眼窝是空的。但在眼窝深处，有一个光点。像是有人在里面点过一根蜡烛又吹灭了。不，不是蜡烛。是反光。黄色的。非常微弱。

\noindent 而在尸体的胸口——用一条腐烂的皮条绑着，勒进骨头里——有一本笔记本。皮条穿过肋骨，在外套底下，像是这个人下葬时书就贴着胸口。我得用刀割断皮条。皮条在我手里碎成了渣。

\noindent 我知道这听起来像谎话。一本1917年随患者下葬、绑在身上的笔记本，六十二年之后，纸页还是\textit{湿润的}？我自己都不会相信。但您正拿着这本笔记本。您看见了纸页。看见了墨迹。

\noindent 您知道我没有撒谎。

\noindent 我拿起笔记本。碰到它的时候，我感到一股热。不是太阳的热。是一种潮湿的热，像揭开伤口的热。仿佛笔记本在里面是活的。在生长。

\noindent 我翻开第一页。

\noindent 笔迹是孩子的。歪歪扭扭，忽上忽下，像自己学着写字的人。顶端有一个签名：

\noindent \textit{若阿金·安东尼奥·科雷亚 —— 患者1,260}

\noindent 还有一个日期。

\noindent \textit{1917 —— 我进这里的那一年}

\noindent 但日记在继续。一页又一页。最后一条记录是1979年的。

\noindent \textit{在我把他挖出来前三天。}

\noindent 您明白吗？这个人已经死了62年。而日记还在被写。在坟里。在黑暗中。用一副骷髅的手指。

\noindent 那一夜我没有睡。

\noindent 直到今天也没有。

\noindent 日记您可以拿走。我不想让它再靠近我。

\noindent --- \textit{热内西奥·巴尔博萨，48岁，巴尔巴塞纳收容医院掘墓人，从业22年。}

\vspace{0.5em}
\noindent\textit{\textbf{↪ 路线:} \rota{MEM-04} \rotasep \rota{LUC-02}}
